\documentclass[10pt,a4paper]{amsart}
\usepackage[margin=19mm]{geometry}
\usepackage{amsmath,amssymb,mathtools}
\usepackage[scr=rsfs,cal=euler]{mathalfa}
\usepackage{xfrac}
\usepackage[unicode,hidelinks,bookmarks=false]{hyperref}
\newcommand{\ie}{i.e.\ }
\newcommand{\cf}{cf.\ }
\newcommand{\resp}{respectively\ }
\newcommand{\Subsection}{\subsection}
\providecommand{\nobreakdash}{\nobreak\hbox{-}}
\setlength{\emergencystretch}{2em}
\multlinegap0pt
\usepackage{algorithm}\usepackage{algpseudocode}
\usepackage{amssymb}
\usepackage{mathtools}
\usepackage{float}
\usepackage{bm}
\usepackage{tikz}
\newtheorem{lemma}{Lemma}
\newtheorem{theorem}{Theorem}
\title{On the computation of the cumulative distribution function of the Normal Inverse Gaussian distribution}
\author{Guillermo Navas-Palencia}
\date{Source: arXiv:2502.16015v2 (CC BY 4.0)}
\begin{document}
\maketitle

\noindent\emph{Source and licence.} On the computation of the cumulative distribution function of the Normal Inverse Gaussian distribution. Version arXiv:2502.16015v2 (CC BY 4.0), distributed by arXiv under the Creative Commons Attribution 4.0 International licence, http://creativecommons.org/licenses/by/4.0/, as recorded in the arXiv licence metadata for this exact version and shown on the abstract page https://arxiv.org/abs/2502.16015v2. Full licence notice: \texttt{licenses/arxiv-2502.16015-CC-BY-4.0.txt}.\par\smallskip

\noindent\textit{Editorial scope.} Counted unit: the article's theorem theorem\_expansion\_xmu\_b0\_postivie\_remainder (source0.tex lines 284--289). The article's complete original proof of it is delivered with it (source0.tex lines 291--310, one \texttt{proof} environment of 20 lines). No mathematics is written, repaired or re-derived: the statement and the proof are the article's own lines, in the article's own notation, and no retained line is substituted or re-flowed.  Retained uncounted as the source-local context the proof needs: lines 258--260; lines 263--268.  Nothing was omitted from any retained span, so the package carries no omission record.  The retained lines cite 1 works, whose entries are transcribed verbatim from the article's own bibliography source in the e-print and delivered with short editorial labels recorded in the item's reference mapping.  No retained line contains a figure environment, an image-inclusion command or drawing code, so no image is delivered and the package declares no asset dependency.  Editorial formatting only, in addition: an AMS \texttt{amsart} preamble that restates the article's own declarations and macros used by the retained lines, namely the environments \texttt{lemma}, \texttt{theorem} on separate counters, together with the standard packages those lines need.  The retained environments print in this document as Lemma 1, Theorem 1 in order of appearance, whereas the article prints the same items with the numbering of its own full document.  The complete native source of the article as distributed in the arXiv e-print for this exact version is bundled unchanged as \texttt{sources/173692/source0.tex}.
\begin{equation}\label{expansion_xmu_b0_positive_with_remainder}
F(x; \alpha, 0, \mu, \delta) = \frac{1}{2} + \frac{\delta e^{\delta \alpha}}{\pi}\frac{(x-\mu)\alpha}{\omega}\left(\sum_{k=0}^{N-1}T_k + R_N\right).
\end{equation}

\begin{lemma}\label{lemma_1}
For $x \ge 0$ and $\nu \ge -\frac{1}{2}$ we have
\begin{equation}
K_{\nu + 1}(x) < \frac{\Gamma(\nu + 1)2^{\nu}}{x^{\nu + 1}}.
\end{equation}
\end{lemma}

\begin{theorem}\label{theorem_expansion_xmu_b0_postivie_remainder}
Given $\alpha > 0$, $\omega > 0$, $x-\mu \in \mathbb{R}$ and $N \in \mathbb{N}$, the first omitted term $T_N$ in \eqref{expansion_xmu_b0_positive_with_remainder} satisfies
\begin{equation}\label{bound_TN_xmu_b0_positive}
T_N < \frac{1}{2\alpha\omega}\left(\frac{x-\mu}{\omega}\right)^{2N} \sqrt{\frac{\pi}{N + 1/2}}.
\end{equation}
\end{theorem}

\begin{proof}
The use of Lemma \ref{lemma_1} gives
\begin{align*}
T_N &= \left(\frac{(x-\mu)^2 \alpha}{\omega}\right)^N \frac{K_{N+1}(\alpha \omega)}{(2N +1)!!}\\
&<  \left(\frac{(x-\mu)^2 \alpha}{\omega}\right)^N \frac{\Gamma(N + 1)2^N}{(\alpha\omega)^{N + 1}(2N +1)!!}\\
&= \frac{1}{\alpha\omega}\left(\frac{x-\mu}{\omega}\right)^{2N} \frac{(N! 2^N)^2}{(2N + 1)!}.
\end{align*}
An upper bound for the ratio of factorials in the previous inequality is given by
\begin{equation*}
\frac{(N! 2^N)^2}{(2N + 1)!} = \frac{\sqrt{\pi}}{2}\frac{\Gamma(N+1)}{\Gamma(N + 3/2)} \le \frac{1}{2}\sqrt{\frac{\pi}{N + 1/2}},
\end{equation*}
where we use the fact that $\Gamma(N + 3/2) = (N + 1/2) \Gamma(N + 1/2)$ and the upper bound of the ratio of gamma functions \cite{Wendel1948}
\begin{equation}
\frac{\Gamma(x + 1)}{\Gamma(x+s)} \le (x + s)^{1-s}, \quad s \in (0, 1).
\end{equation}
Thus, the following bound for $T_N$ holds
\begin{equation}
T_N < \frac{1}{2\alpha\omega}\left(\frac{x-\mu}{\omega}\right)^{2N} \sqrt{\frac{\pi}{N + 1/2}}.
\end{equation}
\end{proof}

\begin{thebibliography}{99}
\bibitem[Wendel]{Wendel1948} J.~G. Wendel. \newblock {Note on the Gamma Function}. \newblock {\em American Mathematical Monthly}, 55:563, 1948.
\end{thebibliography}
\end{document}
